\subsection{Лемма Гейне-Бореля.}

\subsubsection*{Определения}
\begin{tcolorbox}
    \textbf{Открытое множество} $G$: множество, каждая точка которого является внутренней.
    \[ \forall x \in G \;\; \exists U_\varepsilon(x) \subset G \]
    (Точка входит в множество вместе с некоторым "запасом" — окрестностью).
\end{tcolorbox}

\begin{tcolorbox}
    \textbf{Замкнутое множество} $F$: множество, которое содержит все свои предельные точки.
    \[ \text{Если } x_n \in F \text{ и } x_n \to x_0, \text{ то } x_0 \in F \]
    (Из замкнутого множества "нельзя убежать" при переходе к пределу).
\end{tcolorbox}

\begin{tcolorbox}
    \textbf{Покрытие}: Система множеств $\{G_\alpha\}$ называется покрытием $F$, если $F \subset \bigcup_\alpha G_\alpha$.
\end{tcolorbox}

\subsubsection*{Лемма Гейне-Бореля}
\begin{tcolorbox}
    Пусть $F \subset \mathbb{R}$ — \textbf{замкнутое и ограниченное} множество (например, отрезок).
    Пусть $\{G_\alpha\}$ — бесконечное открытое покрытие множества $F$.
    \textbf{Тогда} из него можно выделить \textbf{конечное подпокрытие}.
    \[ \exists G_{\alpha_1}, G_{\alpha_2}, \dots, G_{\alpha_k}: \quad F \subset \bigcup_{i=1}^k G_{\alpha_i} \]
\end{tcolorbox}

\begin{tcolorbox}[title=Доказательство (Метод вложенных отрезков)]
    Доказываем от противного.
    
    1. \textbf{Предположение:} Пусть существует такое открытое покрытие $\{G_\alpha\}$, из которого \textbf{нельзя} выбрать конечное подпокрытие.
    
    2. Так как $F$ ограничено, поместим его в отрезок $[a_0, b_0] \supset F$.
    
    3. \textbf{Шаг деления:} Разобьем отрезок $[a_0, b_0]$ пополам точкой $c_0$.
    Тогда часть множества $F$, лежащая в $[a_0, b_0]$, представится как объединение:
    \[ F \cap [a_0, b_0] = (F \cap [a_0, c_0]) \cup (F \cap [c_0, b_0]) \]
    Если бы обе эти части покрывались конечным числом множеств из $\{G_\alpha\}$, то и всё множество $F$ покрывалось бы их суммой (конечное + конечное = конечное).
    Но мы предположили, что конечного покрытия нет. Значит, хотя бы одна из частей \textbf{не допускает} конечного покрытия.
    Обозначим эту «плохую» половину $[a_1, b_1]$.
    
    4. \textbf{Продолжение процесса:}
    Делим $[a_1, b_1]$ пополам. Выбираем ту половину $[a_2, b_2]$, пересечение которой с $F$ нельзя покрыть конечным числом множеств.
    
    Получаем систему вложенных отрезков:
    \[ [a_0, b_0] \supset [a_1, b_1] \supset \dots \supset [a_n, b_n] \supset \dots \]
    Длина отрезка стремится к нулю: $|b_n - a_n| = \frac{b_0-a_0}{2^n} \to 0$.
    \textbf{Свойство:} Для любого $n$ множество $F \cap [a_n, b_n]$ не покрывается конечным числом $G_\alpha$.
    
    5. \textbf{Предельная точка:}
    По принципу вложенных отрезков (лемма Кантора), существует единственная точка $c$, принадлежащая всем отрезкам:
    \[ c = \bigcap_{n=1}^\infty [a_n, b_n] \]
    
    6. \textbf{Принадлежность $c$ множеству $F$:}
    В каждом отрезке $[a_n, b_n]$ есть точки из $F$ (иначе пересечение было бы пустым и покрывалось бы 0 множеств, что конечно). Значит, $c$ — предельная точка для $F$ (или совпадает с элементами). Так как $F$ \textbf{замкнуто}, оно содержит свои предельные точки $\Rightarrow c \in F$.
    
    7. \textbf{ПРОТИВОРЕЧИЕ (Финал):}
    \begin{itemize}
        \item Так как $c \in F$, а $\{G_\alpha\}$ — покрытие $F$, то существует какое-то открытое множество $G_{\alpha^*} \in \{G_\alpha\}$, которое содержит $c$:
        \[ c \in G_{\alpha^*} \]
        \item Так как $G_{\alpha^*}$ — \textbf{открытое}, то точка $c$ входит в него вместе с некоторой $\varepsilon$-окрестностью:
        \[ (c - \varepsilon, c + \varepsilon) \subset G_{\alpha^*} \]
        \item Так как длина отрезков $[a_n, b_n]$ стремится к нулю и все они содержат $c$, то начиная с некоторого номера $N$ весь отрезок целиком попадет в эту окрестность:
        \[ [a_N, b_N] \subset (c - \varepsilon, c + \varepsilon) \subset G_{\alpha^*} \]
        \item \textbf{Вывод:} Часть множества $F$, попавшая в этот отрезок ($F \cap [a_N, b_N]$), покрывается \textbf{ОДНИМ} множеством $G_{\alpha^*}$.
        \item Один — это конечное число. А мы выбирали этот отрезок так, что его \textbf{нельзя} покрыть конечным числом множеств.
    \end{itemize}
    \textbf{Противоречие.} Значит, исходное предположение неверно, и конечное покрытие существует.
\end{tcolorbox}
